\uuid{2F9q}
\titre{ Étude d'une application linéaire }

\niveau{Licence 1}
\module{ Applications linéaires }
\chapitre{Espaces vectoriels}   			%Continuité, Groupes, Fonctions de plusieurs variables...
\sousChapitre{Matrice et application linéaire}			%Optimisation, Diagonalisation d'une matrice, Calcul de dérivées partielles...

\theme{}				%Fonction de répartition, Division euclidienne de polynômes, ...
\auteur{Quentin Liard}
\datecreate{ {2026-03-26}}
\organisation{}			%AMSCC, Exo7, ...
\difficulte{}			%1, 2, 3, 4 ou 5
\contenu{

\texte{
On considère l’application linéaire $f:\mathbb{R}^3\to\mathbb{R}^3$ définie par
\[
f(x,y,z)=(x+y,y+z,z).
\]
}

\begin{enumerate}

\item \question{Déterminer la matrice de $f$ dans la base canonique.}
\reponse{
\[
A=
\begin{pmatrix}
1&1&0\\
0&1&1\\
0&0&1
\end{pmatrix}.
\]
}

\item \question{Déterminer $\ker f$ et $\mathrm{Im}f$.}
\reponse{$\ker f=\{0\}$ et $\mathrm{Im}f=\mathbb{R}^3$.}

\item \question{$f$ est-elle bijective ?}
\reponse{Oui car $\ker f=\{0\}$.}

\end{enumerate}

}